% Post-CP27 proved input. Four complete actual-file and exact-source audits; genuinely twisted coefficients retained.
\subsection{Second cohomology for arbitrary field characters}

\begin{lemma}[Proved: arbitrary field-character lifting classes]
\label{D:arbitrary-field-character-H2}
Let \(R\) be a fixed complete excellent normal local surface
algebra over an algebraically closed coefficient and residue
field \(k\) of characteristic \(p>0\), and put
\(\Gamma=\pi_1(U_R)\), where
\(U_R=\operatorname{Spec}R-\{\mathfrak m_R\}\).
Let \(H\subset\Gamma\) be open of index \(t\), let
\(q=p^f\), and fix a continuous field character
\[
 \chi:H\longrightarrow\operatorname{Gal}(\mathbb F_q/\mathbb F_p).
                                                        \tag{D.WH.1}
\]
Write \(a=|\operatorname{im}\chi|\), \(\nu=f/a\),
\(q_0=p^\nu\), and \(n=q-1\).
Let \(U_H\) be the puncture corresponding to \(H\), and
let \(V=U_V\to U_H\) be the connected cyclic cover
corresponding to \(\ker\chi\). Its complete normal local
cover ring \(A_V\) has degree \(ta\) over \(R\).
Let \(\mathcal M_\chi\) be the locally constant sheaf
with fibre \(\mathbb F_q^\times\) and its prescribed
field action.

There is a deck-equivariant SNC resolution of \(A_V\)
with \(v_V=O_R(ta)\) exceptional components such that
\[
 \begin{aligned}
 |H^2_{\mathrm{cont}}(H,\mathbb F_q^\times(\chi))|
 &\le |H^2_{\mathrm{et}}(U_H,\mathcal M_\chi)|\\
 &\le(q_0+1)^{\,4p_g(A_V)+3v_V}
 \le q^{c_Rt}.
 \end{aligned}                                         \tag{D.WH.2}
\]
The constant \(c_R\) is independent of \(H,q,\chi,a\).
Neither the cover nor the field character is assumed
tame along an exceptional divisor; in particular
\(p\mid a\) is allowed.
\end{lemma}

\begin{proof}
The case \(n=1\) is immediate. Hence assume \(n>1\).
Fix \(\mathbb F_q^\times=\mu_n(k)\) once, after passing
to \(V\). Choose the geometric deck-pullback convention
\(\tau\) so that the generator of the cyclic descent
group acts on \(H^j(V,\mu_n)\) as
\[
                         \sigma=q_0(\tau^*)^{-1}.
                                                        \tag{D.WH.3}
\]
Thus its invariants are the kernel of \(q_0-\tau^*\).
The geometric automorphism \(\tau\) is \(k\)-linear
and fixes \(\mu_n(k)\). The scalar \(q_0\) in
(D.WH.3) is the coefficient action, not a geometric
action on roots of unity.

\emph{A controlled equivariant SNC model.}
Begin with the minimal regular resolution of \(A_V\).
Every deck automorphism lifts uniquely, by the
minimal-resolution argument in
Lemma~\ref{D:arbitrary-field-character-scalar-H1}.
Consequently these lifts satisfy the group identities,
even when \(p\mid a\).
The all-complete geometric budgets give
\(p_g(A_V)=O_R(ta)\) and \(v_0=O_R(ta)\) for the
number of components on this model.
If \(A_V\) is regular, first blow up its closed point;
this is a single invariant centre and gives a
nonempty proper exceptional curve.

Let \(E_0\) denote the resulting reduced exceptional
curve, and put \(G_0=h^1(E_0,\mathcal O_{E_0})\).
The ideal sequence and coherent \(H^2=0\) give
\(G_0\le p_g(A_V)\). Apply the blowup procedure of
Lemma~\ref{D:linear-SNC-refinement-centres}, at every
step blowing up the entire finite set of non-SNC
points. This set is deck-stable. Its distinct centres
are disjoint, so simultaneous blowup is the same
sequence of individual point blowups in any order.
The proof of that lemma bounds every finite prefix
whose centres are non-SNC, by
\(2\delta_0+7G_0\). It therefore bounds this equivariant
procedure as well, without a factor \(a\).
The resulting SNC model \(Y\) is deck-equivariant,
has \(v=v_V=O_R(ta)\) components and \(O_R(ta)\)
nodes, and has the same geometric genus \(p_g(A_V)\).
No arbitrary pencil or additional orbit of regular
boundary points is used.

Write \(E=\bigcup_{i=1}^v E_i\), \(g_i=g(E_i)\),
and let \(G\) be its connected dual graph, with
\(b_1=b_1(G)\). Put
\(\Lambda=\mathbb Z^{\{E_i\}}\) and
\(N_E=(-(E_i\cdot E_j))_{ij}\).
The deck action on \(\Lambda\) is a permutation
and commutes with \(N_E\).
The exact support and Picard calculation in
Lemma~\ref{D:SNC-general-Brauer-geometry} is natural
for isomorphisms of the regular pair. It gives
the following two equivariant exact sequences:
\[
 \begin{aligned}
 0&\longrightarrow Q_n\longrightarrow H^2(V,\mu_n)
             \longrightarrow B_n\longrightarrow0,\\
 Q_n&=\operatorname{coker}(N_E:\Lambda/n\Lambda
                                      \to\Lambda/n\Lambda),\\
 0&\longrightarrow\prod_i\operatorname{Jac}(E_i)[n](k)
       \longrightarrow B_n
       \longrightarrow H_1(G,\mathbb Z)\otimes\mathbb Z/n
       \longrightarrow0.
 \end{aligned}                                         \tag{D.WH.4}
\]
Here \(B_n=\operatorname{Br}'(V)[n]\).
The constant Tate twist on the last term has been
suppressed using the fixed roots of unity.
This calculation uses proper strict-local base change,
the Kummer sequence, and the divisor/point support
complex of the regular SNC pair. The exact purity
input is \citet[Expose XVI, Theorem 3.1.1 and the
complete proof of Corollary 3.1.3]{A_ILO2014}.
Its Gysin maps and signed incidence map are natural;
no deck-invariant spanning tree is required.
The Jacobian identifications are those of
\citet[Lemma 59.69.1, Tag 03RQ]{StacksLocalPairCohomologyB}.
In particular (D.WH.4) is an equivariant filtration,
not a choice of a splitting of the Brauer sequence.

\emph{Kernel estimates on the equivariant pieces.}
Write \(u=q_0-\tau^*\). For a lattice \(L\) of rank
\(d\) with finite-order automorphism \(T\), Smith
normal form gives
\[
 |\ker(q_0-T:L/nL\to L/nL)|
       \le|\det(q_0-T)|\le(q_0+1)^d.                   \tag{D.WH.5}
\]
Indeed the Smith invariant factors bound the kernel
modulo \(n\), and every complex eigenvalue of \(T\)
is a root of unity.
There is one quotient in (D.WH.4) for which
left exactness alone would not suffice. Set
\(P_n=\Lambda/n\Lambda\) and
\(I_n=\operatorname{im}(N_E\bmod n)\).
Both are \(u\)-stable. The kernel--cokernel sequence
of \(0\to I_n\to P_n\to Q_n\to0\), together with
\(|\operatorname{coker}(u|I_n)|=|\ker(u|I_n)|\), gives
\[
 |\ker(u|Q_n)|
 \le|\ker(u|P_n)|\,|\operatorname{coker}(u|I_n)|
 \le|\ker(u|P_n)|^2
 \le(q_0+1)^{2v}.                                     \tag{D.WH.6}
\]
Thus no estimate for the determinant of the
intersection matrix is introduced.

Let \(J=\prod_i\operatorname{Jac}(E_i)\).
Its deck automorphism has finite order dividing \(a\).
The endomorphism \([q_0]-\tau^*\) is an isogeny:
its product with
\([q_0]^{a-1}+[q_0]^{a-2}\tau^*+\cdots+(\tau^*)^{a-1}\)
is \([q_0^a-1]=[n]\).
The characteristic-polynomial and degree formula give
\[
 |\ker(u|J[n](k))|
 \le\deg([q_0]-\tau^*|J)
 \le(q_0+1)^{2\sum_i g_i}.                            \tag{D.WH.7}
\]
We use \citet[Theorem 10.9 and its written proof,
and Proposition 10.20 with Lemmas 10.21--10.22,
printed pp.46--52]{Milne2008AVScalarB}.
The eigenvalue estimate is on the characteristic-zero
Tate space. No semisimplicity in characteristic
dividing \(a\), or division by \(a\), is used.

For finite \(u\)-modules in an exact sequence
\(0\to A_1\to A\to A_2\to0\), one always has
\(|\ker(u|A)|\le|\ker(u|A_1)|\,|\ker(u|A_2)|\).
Applying this to (D.WH.4), and applying (D.WH.5)
to the integral graph lattice, yields
\[
 |\ker(u:H^2(V,\mu_n)\to H^2(V,\mu_n))|
       \le(q_0+1)^{2v+2\sum_i g_i+b_1}.               \tag{D.WH.8}
\]

We also need a bound on a degree-one cokernel.
The equivariant localization and normalization
filtration used in the proof of
Lemma~\ref{D:arbitrary-field-character-scalar-H1} is
\[
 \begin{gathered}
 0\longrightarrow H^1(E,\mu_n)
 \longrightarrow H^1(V,\mu_n)\longrightarrow J_n'
 \longrightarrow0,\qquad J_n'\subset\Lambda/n\Lambda,\\
 0\longrightarrow H^1(G,\mathbb Z)\otimes\mathbb Z/n
 \longrightarrow H^1(E,\mu_n)\longrightarrow J[n](k)
 \longrightarrow0.
 \end{gathered}                                        \tag{D.WH.9}
\]
For the first line one removes the finite invariant
node set, uses codimension-two purity for \(H^1\),
and then divisor purity for the disjoint smooth
component complements. The last term is a subgroup,
so its \(u\)-kernel is bounded directly by that on
\(\Lambda/n\Lambda\).
Since all these groups are finite, the same kernel
and cokernel orders give
\[
 |\operatorname{coker}(u:H^1(V,\mu_n)\to H^1(V,\mu_n))|
 \le(q_0+1)^{v+2\sum_i g_i+b_1}
 \le(q_0+1)^{v+2p_g(A_V)}.                            \tag{D.WH.10}
\]
The last inequality follows from the normalization
sequence:
\(h^1(E,\mathcal O_E)=\sum_i g_i+b_1\le p_g(A_V)\).
It also shows \(2\sum_i g_i+b_1\le2p_g(A_V)\).

\emph{The genuinely twisted descent calculation.}
Apply Cartan--Leray to the cyclic finite etale cover
\(V\to U_H\), with the sheaf \(\mathcal M_\chi\).
Its exact scheme scope and written proof are
\citet[Theorem 14.9 and Remark 14.10, pp.96--97]
{MilneLEC2013H2D}. The spectral sequence is
\[
 H^i(C_a,H^j(V,\mu_n))\Longrightarrow
                          H^{i+j}(U_H,\mathcal M_\chi),
                                                        \tag{D.WH.11}
\]
where the action is (D.WH.3), not just geometric
pullback. In degree zero the coefficient module is
\(\mathbb F_q^\times\) with action \(q_0\).
For the cyclic resolution, its norm is the power
\(n/(q_0-1)\). The norm image has order \(q_0-1\)
and equals the fixed subgroup. Therefore
\[
                       H^2(C_a,\mathbb F_q^\times)=0.
                                                        \tag{D.WH.12}
\]
This computation remains valid when \(p\mid a\).

For any finite \(C_a\)-module \(A\), the cyclic
resolution identifies
\(H^1(C_a,A)=\ker(\mathrm{Norm})/(\sigma-1)A\),
a subgroup of \(A/(\sigma-1)A\).
Since \(\sigma-1=(\tau^*)^{-1}u\), (D.WH.10)
bounds the total-degree-two term with \((i,j)=(1,1)\).
The \((0,2)\) term is the kernel bounded in
(D.WH.8), and (D.WH.12) makes the \((2,0)\)
term zero. Every limiting term is a subquotient
of its second-page term, so no vanishing of the
remaining differentials is assumed. Multiplying
the two bounds gives
\[
 |H^2(U_H,\mathcal M_\chi)|
 \le(q_0+1)^{\,2v+2\sum_i g_i+b_1+v+2p_g(A_V)}
 \le(q_0+1)^{\,4p_g(A_V)+3v}.                         \tag{D.WH.13}
\]
Since \(p_g(A_V)+v=O_R(ta)\), \(q_0+1\le q_0^2\),
and \(q_0^a=q\), this is at most \(q^{c_Rt}\).
In particular the degree \(a\) of the field cover
has been absorbed by the smaller Frobenius exponent.

Finally the degree-two comparison
\[
 H^2_{\mathrm{cont}}(H,\mathbb F_q^\times(\chi))
       \lhook\joinrel\longrightarrow
                  H^2_{\mathrm{et}}(U_H,\mathcal M_\chi)
                                                        \tag{D.WH.14}
\]
is the same Cartan--Leray injection as in the proof
of Lemma~\ref{D:radial-tame-cyclic-H2}; that argument
for the comparison itself does not use tameness.
Pass first to a finite cover making the coefficient
sheaf constant. Every degree-one class on a finite
stage is a finite torsor and dies after a further
finite etale cover, followed by a Galois closure.
Thus the universal pro-cover has no degree-one
cohomology with these coefficients. The five-term
sequence, or its degree-two filtration, gives
(D.WH.14). This proves all of (D.WH.2).
\end{proof}

\begin{remark}[Period and phase scope]
For a prescribed field character, the scalar obstruction
belongs to the twisted group
\(H^2(H,\mathbb F_q^\times(\chi))\).
Its pullback to \(V\) has constant band \(\mu_n\)
and hence an ordinary geometric Brauer image.
A bounded-geometric-period hypothesis for the
semilinear construction refers to this latter image.
No ordinary \(\mathbf G_m\)-Brauer image of the
nonconstant band on \(U_H\) is assumed.
The lemma counts the full twisted coefficient classes.
It is a cohomological input; the actual short-phase
Morita atlas and its specified-Isom comparison are
separate requirements for a semilinear projective count.
\end{remark}
